\documentclass[11pt]{article}

\usepackage[margin=1in]{geometry}
\usepackage{amsmath,amssymb,amsthm}
\usepackage{booktabs}
\usepackage[colorlinks=true,citecolor=blue,linkcolor=blue,urlcolor=blue]{hyperref}
\usepackage[round,authoryear]{natbib}
\usepackage{times}

\newcommand{\eps}{\varepsilon}
\newcommand{\Real}{\mathbb{R}}

\title{FlowSurv-AFT: An Interpretable Deep Accelerated Failure Time Model\\
with Spline-Flow Residuals\\[4pt]
\large Section 1: Introduction}
\author{}
\date{}

\begin{document}
\maketitle

\section{Introduction}
\label{sec:intro}

\subsection{Background and motivation}
\label{sec:background}

Time-to-event, or \emph{survival}, analysis is the branch of statistics
concerned with modelling the time until an event of interest occurs --
death, disease relapse, mechanical failure, customer churn, loan default --
as a function of covariates. It is distinguished from ordinary regression
by a single structural feature that recurs across every application
domain in which it is used: \emph{censoring}. For a substantial fraction
of the subjects under study, the event has not occurred by the end of the
observation window, so the analyst observes not the event time itself but
only a bound on it. Formally, for subject $i$ with true event time $T_i$
and censoring time $C_i$, one observes the pair $(\tilde T_i, \delta_i) =
\bigl(\min(T_i, C_i),\, \mathbb{1}\{T_i \le C_i\}\bigr)$, and any
statistically valid model must incorporate the censored observations --
those with $\delta_i = 0$ -- through their survival probability
$\Pr(T_i > \tilde T_i \mid x_i)$ rather than discard or naively impute
them. A correct treatment of censoring is therefore not an optional
refinement but the defining constraint of the field.

The second defining feature of survival data is that the object of
inferential interest -- the conditional hazard function $h(t \mid x) =
\lim_{\Delta t \to 0} \Pr(t \le T < t + \Delta t \mid T \ge t, x)/\Delta t$
-- takes empirically diverse shapes across the domains in which survival
models are deployed. A hazard may rise monotonically (mechanical wear-out,
aging-related mortality), fall monotonically (early-life or "infant
mortality" failure), follow a \emph{bathtub} shape combining both --
elevated risk shortly after time zero, a stable low-risk middle period,
and rising risk again at longer horizons, the canonical shape in
reliability engineering and in post-surgical or post-treatment oncology
follow-up -- or be \emph{multimodal}, arising when a single labelled
population is in fact a mixture of sub-populations with materially
different risk profiles that a single unimodal hazard family cannot
represent. Whether a given model \emph{can}, in principle, represent
these shapes, independent of how much data is available to fit it, is a
structural property of the model class, not a matter of estimation
efficiency -- and it is the central axis on which this paper positions
its contribution relative to the existing literature.

Two classical model families have dominated applied survival analysis for
five decades, and each buys its strengths with a cost the other does not
pay. The Cox proportional-hazards model \citep{cox1972} is
semiparametric: it estimates covariate log-hazard-ratios by partial
likelihood without committing to a parametric baseline hazard, which
makes it robust to hazard-shape misspecification for the purpose of
estimating relative effects, but leaves it (i) unable to represent hazard
ratios that change over time without explicit extension, and (ii) without
a smooth conditional density, survival curve, or quantile function -- the
baseline hazard is recoverable, if at all, only as a non-smooth
Breslow-type step function. The accelerated failure time (AFT) model
\citep{wei1992aft} takes the opposite route, specifying a complete
parametric law for $\log T$,
\begin{equation}
  \log T = \beta^\top x + \sigma\eps, \qquad \eps \sim F_0,
  \label{eq:aft-intro}
\end{equation}
for a fixed standard distribution $F_0$. This yields a smooth, fully
specified conditional distribution in one pass -- density, survival,
hazard, and quantiles are all closed-form -- and, distinctively, an
effect measure many practitioners find more directly interpretable than a
hazard ratio: the \emph{acceleration factor} $\mathrm{TR} = \exp\{\beta^\top
(x_a - x_b)\}$, the multiplicative factor by which the \emph{entire}
event-time distribution of profile $x_a$ is stretched or compressed
relative to profile $x_b$. In reliability engineering, where lifetime
predictions feed directly into warranty and maintenance-interval
decisions, "expected time-to-failure is multiplied by $\mathrm{TR}$" is
immediately actionable in a way a hazard ratio is not. The AFT's cost is
the fixed choice of $F_0$: a Weibull residual law gives a hazard that is
monotone for every shape parameter but one, and a log-normal residual law
gives a hazard that is unimodal (hump-shaped); neither family can
represent a bathtub or a multimodal hazard, and no amount of data
resolves this -- it is a property of the functional form, not of the
sample size.

Deep learning has been applied to survival analysis from several
directions over the past decade, and -- as this paper's literature review
documents in detail -- each direction inherits a structural limitation
directly traceable to which classical family it extends. Discriminative
deep models built on the Cox partial likelihood (DeepSurv,
\citealp{katzman2018deepsurv}; Cox-Time, \citealp{kvamme2019coxtime})
retain the proportional-hazards assumption and the non-smooth Breslow
baseline; discrete-time deep models (DeepHit, \citealp{lee2018deephit};
PC-Hazard, \citealp{kvamme2021pchazard}) escape proportionality but
discretize time onto a fixed grid, producing step-function survival
curves whose implied density does not exist at arbitrary $t$ and whose
quantiles are resolved only to the grid's bin width. A more recent line
of deep AFT models (deepAFT, \citealp{norman2024deepaft}; DeepR-AFT,
\citealp{kim2022deepraft}; KAN-AFT, \citealp{francis2025kanaft})
modernizes Eq.~\eqref{eq:aft-intro}'s covariate map $\beta^\top x$ into a
neural network, but -- as we establish in detail in
the companion literature review (Sec.~5) -- does so by abandoning the
likelihood itself, estimating instead by iterative imputation,
inverse-probability weighting, or a rank loss, and consequently
forfeiting the smooth, calibrated conditional distribution that was the
AFT's principal advantage over Cox regression in the first place. A
smaller generative literature -- deep mixture models
\citep{nagpal2022dsm}, conditional normalizing flows
\citep{ausset2021flows}, and concurrent discrete-flow work
\citep{licai2026} -- delivers smooth, flexible densities under an exact
censored likelihood, but, as we show in the companion literature review
(Sec.~5--6), each such model
gives up at least one of exact closed-form tractability, AFT
interpretability, or demonstrated fidelity to non-standard hazard shapes.

\subsection{Problem statement}
\label{sec:problem}

No survival model in the published literature, to our knowledge,
simultaneously provides all of the following five properties.

\begin{enumerate}
  \item[\textbf{P1}] An \textbf{arbitrary, learned event-time
    distribution} -- including bathtub, multimodal, and
    covariate-crossing hazard shapes -- represented by a genuinely
    \textbf{smooth density}, not a step function or a grid-resolved
    approximation.
  \item[\textbf{P2}] \textbf{Exact, closed-form} density $f(t\mid x)$,
    survival $S(t\mid x)$, hazard $h(t\mid x)$, quantile function, and
    one-pass sampling -- with no ordinary-differential-equation solver,
    no time discretization, and no numerical root-finding substituting
    for an analytic inverse.
  \item[\textbf{P3}] The \textbf{interpretable acceleration-factor
    decomposition} of Eq.~\eqref{eq:aft-intro} that makes the AFT the
    preferred effect-communication tool over the hazard ratio in
    clinical and reliability practice.
  \item[\textbf{P4}] Estimation by the \textbf{exact full likelihood}
    under right censoring -- the density contribution for observed
    events and the survival contribution for censored observations --
    with no inverse-probability-of-censoring weighting, no partial
    likelihood, and no iterative imputation standing in for the
    likelihood.
  \item[\textbf{P5}] \textbf{Direct empirical evidence} -- absent, we
    argue, from the entire literature reviewed in
    the companion literature review -- of how such a model
    behaves as censoring proportion, censoring mechanism, and hazard
    shape are varied jointly, assessed with modern calibration
    instruments rather than concordance alone, in light of large-scale
    neutral evidence \citep{burk2024neutral} that concordance gains are
    not, in general, available to be won on standard tabular survival
    data.
\end{enumerate}

The companion literature review (Sec.~2--6) develops this claim in full: the
literature separates into four camps -- classical AFT/Royston--Parmar,
discrete-time deep models, likelihood-free deep AFT models, and
generative/flow-based models -- each of which satisfies at most four of
P1--P5, never all five jointly. The present paper is built to occupy
that gap.

\subsection{Motivating context}
\label{sec:motivating}

This work extends a specific, concrete empirical question left open by
prior work. \citet{opatha2024} conducted a careful simulation study
comparing Cox regression against the Weibull AFT across a factorial
design of sample size, censoring proportion, censoring type (Type~I
administrative censoring versus Type~III random censoring), and hazard
shape, and found that the ranking between the two classical models is
\emph{not} fixed: the AFT dominates specifically at small sample sizes
and under monotone increasing hazards, while Cox regression's robustness
to hazard-shape misspecification becomes relatively more valuable as the
true hazard departs further from either model's assumptions. Their study,
however, was confined to two classical models evaluated by Harrell's
concordance alone. It leaves open the natural next question this paper
takes up: \emph{where do deep, flexible-density survival models sit on
this same map, once evaluated with modern calibration and
distributional-fidelity instruments rather than concordance -- and can a
deep AFT model be constructed that retains the AFT's interpretability
while removing the parametric restriction on its residual law entirely,
without sacrificing the exact likelihood that made the classical AFT a
well-calibrated probabilistic model in the first place?}

\subsection{Contributions}
\label{sec:contributions}

This paper makes the following contributions.

\begin{enumerate}
  \item \textbf{FlowSurv-AFT}, a deep accelerated failure time model in
    which the residual law of Eq.~\eqref{eq:aft-intro} is not fixed to a
    standard parametric family but is instead a \emph{conditional
    rational-quadratic spline (RQS) normalizing flow}
    \citep{durkan2019nsf} applied to a standard logistic base
    distribution. Because the RQS transform is a per-bin monotone
    rational-quadratic map, both it and its inverse are available in
    closed form, together with the log-determinant of the Jacobian
    needed for the change-of-variables density -- so FlowSurv-AFT is the
    first model, to our knowledge, to satisfy P1 and P2 simultaneously:
    an arbitrary learned residual shape with no numerical solver in
    \emph{either} direction. The model is fit by the exact
    right-censored full likelihood (P4), and its explicit
    location-scale decomposition $\mu(x) = \beta^\top x$-generalizing
    neural map, together with a fixed-scale, fixed-location base
    distribution, preserves the acceleration-factor interpretation of
    Eq.~\eqref{eq:aft-intro} exactly (P3): $\mathrm{TR}(x_a, x_b) =
    \exp\{\mu(x_a) - \mu(x_b)\}$ continues to multiply the entire fitted
    event-time distribution, not merely a linear predictor.
  \item \textbf{A principled warm start nesting the classical AFT inside
    the deep model.} The residual flow is parameterized so that, at
    initialization, it is exactly the identity map: FlowSurv-AFT begins
    training as a log-logistic AFT and only departs from it as the data
    support additional residual flexibility, which both stabilizes
    optimization and gives a natural non-inferiority baseline (H3 below)
    against which any added flexibility must be judged.
  \item \textbf{The first systematic hazard-shape-recovery and
    calibration study across censoring regimes} (P5) that we are aware
    of in the flexible-density survival literature: a full factorial
    simulation design spanning multiple hazard-shape scenarios (standard
    Weibull, bathtub, multimodal, crossing, and nonlinear-effect
    regimes), sample sizes, censoring proportions from 0 to 80\%, and
    both administrative (Type~I) and random (Type~III) censoring
    mechanisms, evaluated by true-versus-estimated hazard $L_2$ error
    against ground truth, distributional fidelity (Kolmogorov--Smirnov
    and 1-Wasserstein distance to the true conditional CDF), and modern
    calibration instruments (D-calibration; the integrated calibration
    index) alongside a single, pre-registered concordance variant to
    avoid the metric-selection bias \citet{sonabend2022chacking}
    terms "C-hacking."
  \item \textbf{Real-data validation and an interpretability analysis}
    on five public clinical benchmark datasets, assessing whether
    FlowSurv-AFT's learned acceleration factors are concordant with
    classical AFT time ratios where the classical model is adequate
    (extending trust to the deep model's estimates where they diverge),
    and whether its individual survival distributions are better
    calibrated than discriminative deep baselines under heavy censoring
    -- the regime, per \citet{burk2024neutral}, in which this literature
    is most likely to have a genuine advantage to demonstrate.
  \item \textbf{An open, reproducible benchmark} -- code, the full
    120-cell simulation configuration grid, and the resulting
    meta-dataset of fitted-model metrics -- together with
    evidence-based practical guidance on when the added modelling and
    computational cost of a flexible deep survival model is justified,
    calibrated throughout against the neutral-benchmark finding that,
    on standard tabular survival data, discrimination gains alone
    rarely justify it.
\end{enumerate}

\subsection{Research questions and hypotheses}
\label{sec:rq}

The contributions above are organized around four research questions,
each pre-registered with a corresponding directional hypothesis before
any experiment was run.

\begin{description}
  \item[RQ1 / H1a+H1b (flexibility).] \emph{Can a spline-flow residual
    distribution trained by exact censored likelihood recover bathtub,
    multimodal, and crossing conditional hazard shapes with fidelity
    approaching the true data-generating mechanism?} The answer, from a
    Phase 5 pilot and three follow-up diagnostics -- a reduced-budget
    tuning pass, two capacity increases, and a warm-start ablation, all
    converging on the same result and logged as a pre-registration
    deviation -- turns out to depend on the data-generating mechanism,
    not just the hazard's visual shape. \textbf{H1a
    (confirmatory).} On hazard shapes that are residual-shape or
    location perturbations of a common mechanism -- multimodal (S3),
    nonlinear-location (S5) -- FlowSurv-AFT attains lower hazard-recovery
    error than parametric AFT, DSM, and DeepHit. \textbf{H1b
    (descriptive, not confirmatory).} On hazard shapes generated by a
    mechanism not naturally expressible in AFT location-scale-shape form
    -- a Cox-proportional-hazards-multiplicative bathtub (S2), a discrete
    Weibull-shape regime switch producing crossing hazards (S4) -- the
    covariate effect must be absorbed entirely through the flow's
    shape-conditioning pathway rather than the location shift the model
    is warm-started toward, and FlowSurv-AFT's hazard-recovery error
    relative to DSM is reported descriptively, with no directional claim
    of superiority.
  \item[RQ2 / H2 (calibration $\times$ censoring).] \emph{Does the
    calibration advantage of exact-likelihood training over
    discrete-time or partial-likelihood deep training grow with the
    censoring proportion?} We hypothesize that it does: models trained
    by a proper full likelihood should degrade more gracefully under
    heavy censoring than models trained by discretized or
    weighting-based surrogates, whose instability under heavy censoring
    is discussed in the companion literature review (Sec.~5).
  \item[RQ3 / H3 (no flexibility penalty).] \emph{Does FlowSurv-AFT pay
    a cost, in fit quality, for its added flexibility on data the
    classical model already fits well?} We hypothesize not: on
    standard-shaped (Weibull) data, FlowSurv-AFT should match the
    parametric Weibull AFT within a pre-specified non-inferiority
    margin, since the identity-initialized warm start nests the
    classical model inside the deep one.
  \item[RQ4 / H4 (interpretability).] \emph{Do FlowSurv-AFT's learned
    acceleration factors agree with classical AFT time ratios where the
    classical model is adequate, and what do they reveal where it is
    not?} We hypothesize concordance in the former regime (a
    descriptive, not confirmatory, claim; see the pre-registration) and
    document nonlinear acceleration effects the linear classical model
    cannot represent in the latter.
\end{description}

\subsection{Scope and limitations}
\label{sec:scope}

Three scope decisions are stated explicitly rather than left implicit.
First, this study addresses right censoring (Types~I and~III) as the
primary setting, with interval censoring treated as a secondary
experiment; \emph{dependent} (informative) censoring -- where the
censoring mechanism is not independent of the event time given covariates
-- is out of scope and is flagged as an explicit assumption throughout,
with the deep copula-based approach of \citet{gharari2023copula} noted as
the relevant extension point for future work. Second, the study targets
tabular covariates of moderate dimension, consistent with the classical
and machine-learning survival literature this work is positioned
against; image, text, and other structured covariate modalities, and
competing-risks extensions, are left to future work. Third, and most
consequential for how the paper's claims are worded throughout: in light
of the neutral-benchmark evidence of \citet{burk2024neutral} that
discrimination differences among contemporary survival models on
standard tabular data are not, in general, statistically distinguishable
from Cox regression, this paper's claims to real-data predictive value
are deliberately restricted to calibration and distributional quality
rather than concordance -- a discipline enforced mechanically by the
study's pre-registration of its primary confirmatory contrasts before
any experiment was run.

\subsection{Organization of the paper}
\label{sec:organization}

The companion literature review situates FlowSurv-AFT against the four
methodological camps introduced above in full detail (Sec.~2--6), with
particular attention to the two closest prior works,
\citet{ausset2021flows} and the concurrent \citet{licai2026} (Sec.~6),
and closes with a formal statement of the resulting research gap (Sec.~8).
Subsequent sections of the full paper (methodology, simulation design,
real-data experiments, and results, referenced here but developed in the
companion methodology document) present the model's exact closed-form
derivations, the estimation and optimization protocol, the full factorial
simulation study addressing H1--H3, and the real-data validation and
interpretability analysis addressing H4.

\bibliographystyle{apalike}
\bibliography{references}

\end{document}
